\documentclass{article}
% Swap for the target venue's style file when chosen (neurips/icml/tmlr...).
\usepackage[margin=1in]{geometry}
\usepackage{amsmath,amssymb,amsthm}
\usepackage{graphicx,booktabs,microtype}
\usepackage[numbers]{natbib}
\usepackage{hyperref}

\newtheorem{lemma}{Lemma}
\newtheorem{proposition}{Proposition}
\newtheorem{definition}{Definition}

% Notation
\newcommand{\Vf}{V}                                  % value function
\newcommand{\E}{\mathbb{E}}
\newcommand{\Base}{\mathrm{base}}
\newcommand{\Breg}{D_{\varphi}}
\newcommand{\Spence}{\mathrm{Li}_2}                  % dilogarithm

\title{The Spence Loss: Well-Conditioned Value-Function Learning\\
       for Reward-Guided Diffusion}
\author{David Li-Bland \and \ldots}
\date{}

\begin{document}
\maketitle

\begin{abstract}
Reward-guided diffusion and diffusion fine-tuning are, through
entropy-regularized stochastic control, the problem of learning a value
function $\Vf(x,t)=\log\E_{\Base}[\exp r(X_1)\mid X_t=x]$---a conditional
log-partition function whose gradient is the optimal control. The unbiased
Monte-Carlo target for $\Vf$ lives in exponential space, so squared-error
regression has gradients that explode for over-predictions and, worse, vanish
for under-predictions, where the model is most wrong and gradient clipping
cannot help. We observe that every Bregman divergence has the value function
as its population minimizer, and introduce the \emph{Spence loss}%
\footnote{Named for the Spence (dilogarithm) function $\Spence$ appearing in
its value; in our code and experiments it is the \texttt{quad} loss.}: the
Bregman divergence whose gradient has quadratic tails on both sides, with a
closed-form numerically-stable gradient. We place it in a family with squared
error and the Itakura--Saito divergence. A secondary contribution analyzes why
high-dimensional value learning is signal-sparse and studies on-policy
importance-sampling schemes. We validate on 2-moons, a dimensional-scaling
benchmark designed for clean statistical reads, and GMM-40, where our
log-partition estimate is competitive with published diffusion samplers.
\end{abstract}

% !TEX root = ../main.tex
\section{Introduction}
\label{sec:01_intro}

% TODO

% !TEX root = ../main.tex
\section{Setting: value functions for reward-guided diffusion}
\label{sec:02_setting}

% ---------------------------------------------------------------------------
\subsection{The base process and the tilted target}

Fix a time horizon $t\in[0,1]$ and a base diffusion on $\mathbb{R}^{d}$,
\begin{equation}
  \mathrm{d}X_t = f(X_t,t)\,\mathrm{d}t + \sqrt{2a}\,\mathrm{d}W_t,
  \qquad X_0 = 0,
  \label{eq:base-sde}
\end{equation}
with $a>0$ a scalar diffusion coefficient and $W$ a standard $d$-dimensional
Brownian motion. Write $\mathbb{P}$ for the induced path measure, $\E_{\Base}$
and $p_{\Base}$ for expectations and marginals under it, and
$\mathcal{L}=f\cdot\nabla + a\Delta$ for its generator. Given a terminal reward
$r:\mathbb{R}^{d}\to\mathbb{R}$, the object we wish to sample from is the
\emph{reward-tilted} path measure
\begin{equation}
  \frac{\mathrm{d}\mathbb{P}^{\star}}{\mathrm{d}\mathbb{P}}
  \;=\; \frac{e^{r(X_1)}}{Z},
  \qquad
  Z \;=\; \E_{\Base}\!\big[e^{r(X_1)}\big],
  \label{eq:tilt}
\end{equation}
whose terminal marginal is $p^{\star}(x)\propto p_{\Base}(x)\,e^{r(x)}$. This is
the common form of reward-guided generation and of diffusion fine-tuning
\citep{song2021sde,black2024ddpo,uehara2024finetuning}, and it also covers Boltzmann sampling: to target
$\pi\propto e^{-E}$ one sets $r = -E - \log p_{\Base}$, which is how we
instantiate the GMM-40 benchmark in \S\ref{sec:07_experiments}.

Nothing below depends on $f$ having a particular form: if $\mathbb{P}$ is a
pretrained diffusion or flow model, $f$ is its learned drift and the whole
development applies unchanged.\footnote{Note the time convention: $t$ runs from
a deterministic start $X_0=0$ to the model's samples at $t=1$, as in the
diffusion-sampler literature \citep{zhang2021pis,vargas2023dds}, rather than
the denoising direction of DDPM-style presentations.}

\paragraph{The base family we use.}
\label{subsec:base-family}
Our base processes are \emph{pinned interpolants}
\citep{albergo2023interpolants,lipman2023flowmatching}: $X_0=0$, a prescribed
terminal marginal $\nu$, and a Brownian bridge between the two given $X_1=x_1$.
The induced drift $f(x,t)=\E[(X_1-X_t)/(1-t)\mid X_t=x]$ vanishes exactly when
$\nu=\mathcal{N}(0,2aI)$; taking $\nu$ to be a Gaussian mixture instead gives a
base model with nontrivial but analytically known $f$, $\Vf$ and $\log Z$,
which is what makes the study of \S\ref{sec:07_experiments} possible. Details
in App.~\ref{sec:C_experimental_details}.

% ---------------------------------------------------------------------------
\subsection{Control formulation and the value function}

Realizing \eqref{eq:tilt} means adding a drift. For a control $u$ giving the
controlled dynamics $\mathrm{d}X_t=(f+u)\,\mathrm{d}t+\sqrt{2a}\,\mathrm{d}W_t$
with path measure $\mathbb{P}^{u}$, Girsanov's theorem gives
$\mathrm{KL}(\mathbb{P}^{u}\,\|\,\mathbb{P})=\E_{\mathbb{P}^u}\!\int_0^1
\tfrac{1}{4a}\|u_t\|^{2}\,\mathrm{d}t$, so \eqref{eq:tilt} is the solution of
the entropy-regularized control problem \citep{daipra1991stochastic,leonard2014survey}
\begin{equation}
  \max_{u}\ \ \E_{\mathbb{P}^{u}}\!\left[\,r(X_1)
    - \int_0^1 \frac{\|u_t\|^{2}}{4a}\,\mathrm{d}t\right],
  \label{eq:control}
\end{equation}
whose optimum is $\log Z$. The central object is the \emph{soft value function}, the continuous-time
counterpart of the soft value of maximum-entropy reinforcement learning
\citep{haarnoja2017softq,levine2018rlasinference},
\begin{equation}
  \Vf(x,t) \;=\; \log \E_{\Base}\!\big[\,e^{r(X_1)} \,\big|\, X_t=x\,\big],
  \label{eq:value}
\end{equation}
a conditional log-partition function. It solves the HJB equation
$\partial_t \Vf + f\cdot\nabla \Vf + a\Delta \Vf + a\|\nabla \Vf\|^{2}=0$ with
terminal condition $\Vf(\cdot,1)=r$, the optimal control is
\begin{equation}
  u^{\star}(x,t) \;=\; 2a\,\nabla_x \Vf(x,t),
  \label{eq:optimal-control}
\end{equation}
and, since $X_0=0$ is deterministic, $\Vf(0,0)=\log Z$ --- the normalizing
constant, which is the headline metric on external sampler benchmarks
(\S\ref{sec:07_experiments}). Learning $\Vf$ therefore solves the sampling
problem, the fine-tuning problem, and the free-energy-estimation problem at
once, and everything downstream reduces to \emph{regressing $\Vf$ from
Monte-Carlo targets}.

% ---------------------------------------------------------------------------
\subsection{The H-martingale property}
\label{subsec:hmartingale}

Exponentiating \eqref{eq:value} linearizes the HJB equation --- the Hopf--Cole
change of variables underlying linearly-solvable and path-integral control
\citep{kappen2005path,todorov2009efficient}. Writing
$h = e^{\Vf}$, Feynman--Kac gives $\partial_t h + \mathcal{L}h = 0$: $h$ is
space-time harmonic for the base process, so $h(X_t,t)$ is a $\mathbb{P}$-martingale,
and \eqref{eq:optimal-control} is exactly the Doob $h$-transform drift
$f + 2a\nabla\log h$. Concretely, for any $0\le t<s\le 1$,
\begin{equation}
  \boxed{\ e^{\Vf(x,t)} \;=\; \E_{\Base}\!\big[\,e^{\Vf(X_s,\,s)}\,\big|\,X_t=x\,\big]\ }
  \label{eq:hmart}
\end{equation}
with the terminal case $\Vf(\cdot,1)=r$. We call \eqref{eq:hmart} the
H-martingale property; it is the source of every target in this paper, and two
of its features drive the design.

\paragraph{(i) Targets are unbiased in exponential space --- and only there.}
Equation \eqref{eq:hmart} says that a \emph{single} draw of $X_s$ from the base
kernel supplies a scalar $q = \Vf(X_s,s)$ (or $q=r(X_1)$ when $s=1$) satisfying
\begin{equation}
  \E\big[e^{q}\,\big|\,X_t=x\big] \;=\; e^{\Vf(x,t)}.
  \label{eq:unbiased-target}
\end{equation}
The estimator is unbiased for $e^{\Vf}$, not for $\Vf$: the natural log-space
plug-ins are both biased low by a Jensen gap (\S\ref{sec:03_expmse}). This is a
property of the log-partition function rather than of a particular scheme, and
it is what forces the regression into exponential space --- and hence into the
conditioning problem of \S\ref{sec:03_expmse}--\S\ref{sec:04_bregman}.

\paragraph{(ii) The identity is twist-independent.}
Suppose we do not want to draw $X_s$ from the base kernel --- for instance
because base samples rarely reach high-reward regions: importance sampling for
$\E[e^{r}]$ needs on the order of $e^{\mathrm{KL}(\mathbb{P}^{\star}\|\mathbb{P})}$
samples \citep{chatterjee2018sample}, which is the quantitative form of the
sparsity argument of \S\ref{sec:06_sparsity_onpolicy}. Propose instead under any adapted twist
$u$, giving the drift $f+u$ and path measure $\mathbb{P}^{u}$, and correct by
the Radon--Nikodym derivative of the base law with respect to it. By Girsanov's
theorem, restricted to the interval $[t,s]$,
\begin{equation}
  \rho_{t\to s} \;=\; \frac{\mathrm{d}\mathbb{P}}{\mathrm{d}\mathbb{P}^{u}}
  \bigg|_{[t,s]}
  \;=\; \exp\!\left(
    -\frac{1}{2a}\int_t^s u_\tau\cdot \mathrm{d}B_\tau
    \;-\; \frac{1}{4a}\int_t^s \|u_\tau\|^{2}\,\mathrm{d}\tau
  \right),
  \label{eq:girsanov-rho}
\end{equation}
where $\mathrm{d}B_\tau=\sqrt{2a}\,\mathrm{d}W^{u}_\tau$ is the driving noise
under $\mathbb{P}^{u}$ (so the stochastic integral is an It\^o integral against
the noise actually realized by the sampler). Under Novikov's condition
$\rho_{t\to s}$ is an exponential martingale with unit mean, and \eqref{eq:hmart}
becomes
\begin{equation}
  e^{\Vf(x,t)} \;=\; \E_{\mathbb{P}^{u}}\!\big[\,\rho_{t\to s}\, e^{\Vf(X_s,s)}\,\big|\,X_t=x\,\big].
  \label{eq:girsanov}
\end{equation}
In an Euler--Maruyama discretization \eqref{eq:girsanov-rho} factorizes over
steps, each contributing
$-\,u\cdot\Delta B/(2a) - \|u\|^{2}\Delta t/(4a)$ with
$\Delta B\sim\mathcal{N}(0,2a\Delta t\,I)$; this per-step ratio is the
incremental weight of twisted-SMC samplers, a construction that predates its
diffusion applications
\citep{whiteley2014twisted,guarniero2017iapf,heng2020csmc,singhal2025fksteering}. The twist $u$ is therefore \emph{free}: it
changes the variance of the target, never its mean. The proposal is a
variance-reduction device, not a modeling choice --- which is what licenses the
importance-sampling and SMC machinery of \S\ref{sec:06_sparsity_onpolicy}
without changing the object being learned.

% ---------------------------------------------------------------------------
\subsection{How training pairs are formed}
\label{subsec:targets}

A learner $\Vf_\theta$ is fit to pairs $(x_t, q)$, of two kinds. \emph{Anchored} targets terminate at the true reward and satisfy
\eqref{eq:unbiased-target} exactly, so $e^{q}$ is unbiased for $e^{\Vf(x_t,t)}$
however wrong the current $\Vf_\theta$ is. \emph{Bootstrapped} targets
substitute the learner's own value at a later time, and satisfy no such
identity. Both appear below; the distinction matters, and we are explicit about
which guarantee is in force.

\paragraph{Off-policy (bridge) anchors.}
Draw a terminal point $x_1\sim\nu$ from the base terminal marginal, place $x_t$
on the base bridge to it, and take $q=r(x_1)$. For the pinned interpolants of
\S\ref{sec:02_setting} that bridge is the Brownian bridge in closed form,
\begin{equation}
  x_t \;=\; t\,x_1 + \sqrt{2a\,t(1-t)}\;\varepsilon,
  \qquad \varepsilon\sim\mathcal{N}(0,I),
  \label{eq:fwd-bridge}
\end{equation}
so the sampled joint $(x_t,x_1)$ is \emph{exactly} the base joint and $q$ is
anchored in the sense above. What matters is that $x_1$ is drawn from the base
marginal $\nu$: anchoring on a reward-aware or held-out data distribution
instead breaks the conditional and biases the targets.

These anchors are cheap \emph{per sample} --- no trajectory is integrated, and
each pair costs a single evaluation of $r$ --- and they cover the whole
marginal. Two caveats. First, $x_1$ is drawn without regard to $r$, which is
what fails in high dimension (\S\ref{sec:06_sparsity_onpolicy}). Second, cheap
per sample is not the same as cheap in total: because the anchors are
uninformative about where $r$ is large, reaching a given accuracy can take many
more of them, and hence many more reward evaluations, than a targeted scheme.
When $r$ is itself expensive --- a quantum-chemistry calculation, a docking
score, a simulator rollout --- total calls to $r$, not wall-clock per sample,
is the budget that binds, and the ranking between the constructions can invert.
Our benchmarks all have cheap analytic rewards, so we report sample counts
rather than reward-call counts; we flag this as a limitation in
\S\ref{sec:09_conclusion}.

\paragraph{On-policy (bootstrapped) targets, and backward-noising.}
The other two constructions are used by the experiments but are not needed to
follow the argument, so we state them briefly and defer the details to
App.~\ref{sec:B_onpolicy_zoo}. A twisted SMC sweep supplies \emph{bootstrapped}
targets, formed by aggregating the learner's own values at the next step; these
satisfy no exact identity, being unbiased in exponential space for the one-step
backup of the current $\Vf_\theta$ rather than for $\Vf$, and their usefulness
rests on the terminal step carrying $q=r(X_1)$ exactly. Separately, because
$e^{\Vf}$ is harmonic, a target attached at $(x_s,s)$ can be reattached at a
point noised backwards from it through the base backward kernel
\begin{equation}
  X_t \,\big|\, X_s = x_s \;\sim\; \mathcal{N}\!\Big(\tfrac{t}{s}\,x_s,\;
    2a\,\tfrac{t(s-t)}{s}\,I\Big), \qquad t<s,
  \label{eq:backward-kernel}
\end{equation}
which for the pinned interpolants does not depend on the terminal law $\nu$ at
all (App.~\ref{sec:D_proofs}), so the construction is available for base models
with nontrivial drift. It is exact when the noised point is base-distributed
and approximate under a twist (\S\ref{sec:06_sparsity_onpolicy}).

\medskip
\noindent In every case the learning problem has the same shape: given pairs
$(x_t,q)$ with $\E[e^{q}\mid x_t] = e^{\Vf(x_t,t)}$, fit $p=\Vf_\theta(x_t)$.
The rest of the paper concerns the two remaining degrees of freedom --- the
divergence applied to $(p,q)$, and the sampler that produces the pairs.

% !TEX root = ../main.tex
\section{The problem with exponential-space squared error}
\label{sec:03_expmse}

% TODO

% !TEX root = ../main.tex
\section{Bregman divergences for value learning}
\label{sec:04_bregman}

% ---------------------------------------------------------------------------
\subsection{Any Bregman divergence learns the value function}

For a strictly convex, differentiable potential $\varphi$, the Bregman
divergence is
\begin{equation}
  \Breg(u, v) = \varphi(u) - \varphi(v) - \varphi'(v)\,(u - v).
\end{equation}
Squared error is the special case $\varphi(u)=u^2$. The property we exploit:

\begin{lemma}[Bregman divergences are proper losses for the value]
\label{lem:minimizer}
Let the regression target be $e^{q}$ with $q$ a conditionally-unbiased
estimate of the value (the H-martingale target, \S\ref{sec:02_setting}). For
any strictly convex $\varphi$,
\begin{equation}
  \arg\min_{u}\ \E_{q\mid x_t}\!\big[\Breg(u,\, e^{q})\big] \;=\; \E[e^{q}\mid x_t],
\end{equation}
so the minimizing prediction in log space is $p^\star=\log\E[e^{q}\mid x_t]=\Vf(x_t)$.
\end{lemma}
\noindent Thus the choice of $\varphi$ does \emph{not} bias the learned value;
it is free to be chosen for \emph{conditioning}. (Proof in
App.~\ref{sec:D_proofs}.)

% ---------------------------------------------------------------------------
\subsection{Design criteria and two reference points}

We parameterize the prediction in log space, $u=e^{p}$, and ask how the
gradient $\partial L/\partial p$ behaves as the log-error $p-q\to\pm\infty$.

\paragraph{Squared error ($\varphi(u)=u^2$).}
$\partial L/\partial p = 2e^{p}(e^{p}-e^{q}) = 2e^{2p}-2e^{p+q}$, which
\emph{explodes} as $p\to+\infty$ and \emph{vanishes} as $p\to-\infty$ --- the
under-prediction pathology of \S\ref{sec:03_expmse}.

\paragraph{Itakura--Saito ($\varphi(u)=-\ln u$).}
$D_{\mathrm{IS}}(e^{q},e^{p}) = e^{q-p}-(q-p)-1$; its gradient in $p$ is
$1-e^{q-p}$, a \emph{linear} (bounded-gradient) tail for over-predictions that
removes the vanishing problem, but it still explodes for under-predictions.
It is \emph{scale-invariant} ($D_{\mathrm{IS}}(cu,cv)=D_{\mathrm{IS}}(u,v)$)
and is the negative log-likelihood of a gamma/exponential observation model,
giving a natural uncertainty interpretation for $\Vf$. A useful midpoint, but
one-sided.

% ---------------------------------------------------------------------------
\subsection{The Spence loss}
\label{subsec:spence}

We propose the Bregman divergence generated by the potential $F$ with
\begin{equation}
  F''(u) = -\,\frac{\ln u}{u\,(1-u)},
\end{equation}
applied to $u=e^{p}$ against the target $e^{q}$. Its gradient with respect to
the prediction has the closed form
\begin{equation}
  \boxed{\ \frac{\partial L}{\partial p} \;=\; p\,\frac{e^{p}-e^{q}}{e^{p}-1}\ }
  \label{eq:spence-grad}
\end{equation}
which is zero iff $p=q$ and, by Lemma~\ref{lem:minimizer}, has population
minimizer $\Vf$.

\paragraph{Two-sided quadratic tails.}
From \eqref{eq:spence-grad},
\begin{align}
  p\to+\infty:\quad & \frac{\partial L}{\partial p}\to p
    &&\text{(no $e^{2p}$ blow-up; over-prediction is quadratically penalized),}\\
  p\to-\infty:\quad & \frac{\partial L}{\partial p}\to p\,e^{q}
    &&\text{(non-vanishing corrective signal for under-prediction).}
\end{align}
Both tails are linear in $p$, i.e.\ the loss is quadratic on both sides --- the
first loss to fix both pathologies of exponential-space regression at once.

\paragraph{Value via the Spence function.}
The divergence value (used only for monitoring; the gradient uses
\eqref{eq:spence-grad}) is expressible through the Spence function
$S(x)=\Spence(1-e^{x})$, using the dilogarithm identity
$S'(x)=-\,x e^{x}/(e^{x}-1)=x/\mathrm{expm1}(-x)$; hence the name. Full
derivation, the numerically stable branch-wise gradient (finite for all float
inputs, unlike the naive autograd which NaNs for $p\gtrsim 88$), and the
gamma-like uncertainty interpretation are in App.~\ref{sec:A_spence}.

% ---------------------------------------------------------------------------
\begin{table}[t]
\centering
\caption{The value-learning loss family. Tails give $\partial L/\partial p$ as
the log-error $\to\pm\infty$. All three share the same minimizer,
$\Vf=\log\E[e^{q}\mid x_t]$ (Lemma~\ref{lem:minimizer}).}
\label{tab:loss-family}
\begin{tabular}{lcccc}
\toprule
loss & under-predict ($p\!\to\!-\infty$) & over-predict ($p\!\to\!+\infty$)
     & scale-inv. & uncertainty model \\
\midrule
squared error   & vanishes ($e^{p+q}$)     & explodes ($e^{2p}$) & no  & Gaussian (exp) \\
Itakura--Saito  & explodes                 & linear (bounded)    & yes & gamma \\
\textbf{Spence (ours)} & \textbf{quadratic} ($p\,e^{q}$) & \textbf{quadratic} ($p$) & no & gamma-like \\
\bottomrule
\end{tabular}
\end{table}

% !TEX root = ../main.tex
\section{The loss in isolation}
\label{sec:05_offpolicy_exp}

\S\ref{sec:04_bregman} is an argument about conditioning, and conditioning
arguments are cheap to make and easy to get wrong. This section measures the
five losses against each other directly.

\paragraph{Design.}
Every arm trains the same network on the same problem instances from the same
off-policy bridge anchors with the same budget; only \texttt{loss\_type}
differs. Training is off-policy throughout, deliberately: an on-policy sampler
would let the loss change the data it is subsequently trained on, and we want
the divergence isolated from the sampler. The arms are the Spence loss,
exponential-space squared error, Itakura--Saito, generalised KL
($\beta=1$), and log-space squared error.

\emph{No arm uses gradient clipping.} This is deliberate and it is what makes
the comparison a test of the tails rather than of a clipping heuristic:
Itakura--Saito and generalised KL each explode on one side
(\S\ref{sec:04_bregman}), and clipping is the obvious remedy for both. Whether
a clipped competitor closes the gap is the natural follow-up and we have not
run it, so nothing here should be read as a claim against clipped variants.

Two of these choices deserve comment. The squared-error arm is our
\emph{stabilised} implementation --- exact gradient, saturated to the float
range --- and not the naive one. Comparing against the naive version would
measure arithmetic rather than the divergence, and would flatter us: the naive
loss simply fails. The log-space arm is included as a bias probe rather than as
a candidate. It is perfectly conditioned and converges happily; by
\eqref{eq:jensen-gap} it converges to the wrong function, and the point of
including it is to show that deficit rather than to propose it.

\paragraph{Tuning, and why it matters here.}
The shared hyperparameters of the benchmark were originally fitted with the
Spence loss in place. Evaluating the alternatives at those values would be a
rigged comparison, and it is the first thing a reader should suspect. We
therefore re-tune the learning rate per (loss, dimension) over a grid spanning
$10^{-5}$ to $3\times10^{-3}$ on separate tuning seeds at a shorter budget, and
run the reported grid on fresh seeds at the chosen value. Each arm is thus
reported at its own best learning rate rather than at ours.

One caveat about that grid. At the larger dimensions several arms --- including
ours --- select the smallest rate offered, $10^{-5}$, so their optima may lie
outside the range and the tuning is truncated. It is truncated identically for
every arm, so the comparison remains matched, and it is a comparison of best
rates \emph{within} $[10^{-5},3\times10^{-3}]$ rather than of global optima. We
mention it because a boundary optimum is exactly the kind of detail that makes
a tuned-versus-tuned claim less clean than it looks.

\paragraph{Metrics.}
Two axes, because they can disagree, and the reason they can matters. The
control only ever sees $\nabla \Vf_\theta$, so a value function that is wrong
by a slowly-varying offset still steers correctly; conversely a loss can score
well on control while learning a badly biased value. The pathology of
\S\ref{sec:03_expmse} concerns the \emph{magnitude} of the learned value in
under-predicted regions, so the value axis is where it should show, and a
control-only comparison could miss it entirely.

\emph{Control quality} is the fraction of the calibrated $6$-nat headroom
closed, the same score used throughout \S\ref{sec:07_experiments}. \emph{Value quality} is the RMSE of $\Vf_\theta$
against the \emph{analytic} value over the base marginal --- available because
this problem family has $\Vf$ in closed form (App.~\ref{sec:C_experimental_details}) ---
together with its signed mean, which exposes bias rather than variance. We also
count non-finite batches, the stability axis of \S\ref{sec:03_expmse}.

% !TEX root = ../main.tex
% Auto-generated placeholder; replaced by make_ablation_table.py once the
% loss-ablation grid completes.
\begin{center}\itshape
[Results table pending: the loss-ablation grid is still running.]
\end{center}


\paragraph{What the comparison shows.}
Four readings, all paired on the same instances. The first is the one we set
out to test; the third is not the one we expected.

\emph{The Spence loss reliably repairs exponential-space squared error.} On
value RMSE it is better at every dimension we ran ($-0.37$ to $-0.75$ nats,
significant at $d=2$ and $d=32$), and on control it is ahead at every
dimension. This is the claim of \S\ref{subsec:spence} and it holds: a weight
with both tails controlled learns a more faithful value than one whose
under-prediction gradient vanishes.

\emph{It also beats generalised KL on value, everywhere.} The $\beta=1$ member
is tied with ours on control at every dimension, but its value RMSE is worse at
all four ($-0.10$ to $-0.23$ in our favour). Fixing one tail is not the same as
fixing both, and the difference shows up where the theory says it should.

\emph{But it does not dominate the family: Itakura--Saito is the better value
estimator above $d=2$.} We win that comparison at $d=2$ ($-0.12$,
significant) and lose it at $d=8$, $32$ and $128$ ($+0.11$, $+0.20$, $+0.16$;
significant at $d=128$), with control statistically tied throughout. The
direction is consistent across three dimensions, so we do not read it as
noise. A reader is entitled to conclude that on this benchmark, above two
dimensions, a divergence from 1968 estimates the value at least as well as
ours. What remains true and is worth separating out: Itakura--Saito achieves
this with a gradient that explodes for under-predictions, so it depends on the
clipping or saturation that our weight makes unnecessary; we did not clip any
arm (App.~\ref{subsec:a-numerics} saturates only where the true gradient
leaves the float range). Whether the tail control buys anything once a
competitor is properly clipped is the experiment we did not run, and it is the
one we would run next.

\emph{Conditioning and unbiasedness are different goods.} The log-space arm is
the strongest controller in the table at every dimension --- by $12$ to $26$
points --- while carrying a value bias of $1.6$ to $4.1$ nats in the direction
\eqref{eq:jensen-gap} predicts. A loss can be well conditioned and wrong. If
the value is wanted only as a gradient that trade is available and cheap; if
$\Vf$ itself is the object --- and $\Vf(0,0)=\log Z$ --- it is not a trade at
all, which is the argument of \S\ref{sec:03_expmse} made empirical.

\paragraph{One axis separates nothing.}
We counted non-finite batches per arm and found none, in any arm, at any
dimension. Once the stabilised implementations of \S\ref{sec:03_expmse} and
App.~\ref{subsec:a-numerics} are used, the overflow failures that motivated
this work do not recur, and the differences above are differences in what the
losses \emph{learn}, not in whether they run. The stability axis motivates the
numerics and is then removed by them.

\paragraph{Scope.}
This experiment isolates the divergence by construction, and therefore says
nothing about how the losses interact with an on-policy sampler, where the
loss influences its own future training data. It is also a single problem
family at one difficulty setting. The dimensional trend within it is
informative; extrapolating it to other families is not warranted.

% !TEX root = ../main.tex
\section{Signal sparsity and on-policy learning}
\label{sec:06_sparsity_onpolicy}

% TODO

% !TEX root = ../main.tex
\section{Experiments}
\label{sec:07_experiments}

% TODO

% !TEX root = ../main.tex
\section{Related work}
\label{sec:08_related}

\paragraph{Reward-guided diffusion as stochastic control.}
Casting diffusion fine-tuning as entropy-regularized control, with the tilted
path measure \eqref{eq:tilt} as the optimum, follows
\citet{uehara2024finetuning}; the underlying variational identity is classical
\citep{daipra1991stochastic,leonard2014survey}. That the exponentiated value is
space-time harmonic, so that the optimal drift is a Doob $h$-transform, is the
Hopf--Cole structure of linearly-solvable and path-integral control
\citep{kappen2005path,todorov2009efficient}, and $\Vf$ is the continuous-time
soft value of maximum-entropy reinforcement learning
\citep{haarnoja2017softq,levine2018rlasinference}. Our contribution is not the
setting but what to do with the regression problem it hands you.

\paragraph{Losses for diffusion samplers.}
The closest competing work attacks the conditioning of value learning by
changing the \emph{objective} rather than the divergence. Log-variance and
adjoint losses \citep{nusken2021logvariance,richter2020vargrad,domingoenrich2024adjoint,domingoenrich2023socm}
replace the pointwise regression with a path-space functional; the
trajectory-balance family
\citep{malkin2022trajectory,madan2023subtb,venkatraman2024rtb} does the same
through a flow-consistency constraint. \citet{zhang2024dgfs} is nearest of all:
it learns intermediate log-partition functions for a diffusion sampler, which
is our object, via partial-trajectory optimization. These are well-motivated
and in several cases better than what we do on sample quality. The distinction
is narrow and worth stating plainly: they buy conditioning by giving up the
conditionally-unbiased exponential-space target, whereas we keep that target
and change only the divergence applied to it. Which trade is better is an
empirical question we do not settle here; what we show is that the second
option exists and is cheap, since it is a drop-in replacement for a loss
function.

\paragraph{Badly-scaled value regression.}
That regressing value targets spanning many orders of magnitude is unstable is
known in reinforcement learning, and the standard remedies transform the
target: adaptive normalization \citep{vanhasselt2016popart}, invertible
rescaling \citep{pohlen2018observe}, or replacing regression with
classification over bins \citep{farebrother2024stopregressing}. All three work
well and none is available here, for the reason given in
\S\ref{sec:03_expmse}: each destroys unbiasedness of the exponential-space
target, which is the only thing making the minimizer the value.

\paragraph{Bregman divergences.}
The divergence is \citet{bregman1967relaxation}; that its expectation is
minimized by the mean, and that Bregman divergences are exactly the losses
with this property, is the Savage representation of proper scoring rules
\citep{savage1971elicitation,gneiting2007proper,banerjee2005clustering}. The
correspondence with exponential families is
\citet{banerjee2005clustering,azoury2001relative}. Itakura--Saito
\citep{itakura1968analysis} and the $\beta$-divergence family
\citep{fevotte2009isnmf,fevotte2011beta,cichocki2010families} are the standard
alternatives; \S\ref{sec:04_bregman} shows why no member of that family has the
tail behaviour we need. We are not aware of prior use of the potential
\eqref{eq:a-phi} as a loss.

\paragraph{Twisted SMC.}
Using a learned twist to steer a particle system, and correcting with the
per-step kernel ratio, is the twisted particle filter
\citep{whiteley2014twisted,guarniero2017iapf} and controlled SMC
\citep{heng2020csmc}; the diffusion-model instantiations are
\citet{wu2023tds,zhao2024twisted,singhal2025fksteering,skreta2025fkcorrectors},
and \citet{wang2026tritsmc} adds the trust-region machinery we evaluate in
App.~\ref{sec:B_onpolicy_zoo}. Our use is standard; the contribution there is
the coverage-versus-concentration finding, which runs against the natural
intuition that a sharper twist should help.

\paragraph{Diffusion samplers.}
The benchmark targets and baselines come from the neural-sampler literature
\citep{zhang2021pis,vargas2023dds,richter2023dis,midgley2021fab,akhoundsadegh2024idem,sendera2024offpolicy,chen2024scld},
evaluated following the protocol concerns raised by
\citet{blessing2024beyondelbos}.

% !TEX root = ../main.tex
\section{Limitations and conclusion}
\label{sec:09_conclusion}

The value function of a reward-guided diffusion is a conditional log-partition
function, and the only unbiased Monte-Carlo targets for it live in exponential
space. That fact is what forces the awkward regression, and squared error
handles it badly in a specific and asymmetric way: over-predictions explode,
which clipping survives, and under-predictions vanish, which nothing survives.
Since every Bregman divergence has the value as its minimizer, the potential is
free, and in log space the entire family reduces to one positive weight on the
residual. Choosing that weight to be $\ln y/(y-1)$ gives gradients linear in the
prediction on both sides, in closed form, with a numerically stable
implementation. It is a drop-in replacement for a loss function.

What the measurements support is narrower than that construction might suggest,
and we would rather state it than let the framing carry it. The Spence weight
repairs exponential-space squared error on both metrics at every dimension we
ran, and beats generalised KL on value fidelity, so controlling both tails is
demonstrably better than controlling one. It does not dominate the family. We ran the
experiment that would decide this --- the same grid with gradient clipping on
every arm --- and it decided against us: clipped Itakura--Saito is the better
controller above $d=2$ and ties on value, and clipping turns out to help our
loss substantially too, which a weight that genuinely needed no clipping should
not do. The durable contribution is therefore the reduction itself --- that
loss design here is the choice of a single positive weight, and that squared
error is a poor choice of it --- and not a claim that our weight is the best
available one. We would recommend the reduction to a practitioner, and clipped
Itakura--Saito as the loss.

\paragraph{Limitations.}
The empirical claims are narrower than the analysis.

\emph{Reward-call budget.} All our targets have cheap analytic rewards, so we
measure cost in samples seen. Under an expensive reward --- a quantum-chemistry
calculation, a docking score, a simulator rollout --- the binding budget is
total evaluations of $r$, and the ranking between off-policy anchors and
on-policy sampling could differ from what we report
(\S\ref{subsec:targets}).

\emph{Transfer.} The on-policy recipe was fitted on one benchmark family, and
does not transfer unchanged to the Boltzmann targets of
\S\ref{subsec:gmm40}. We report the fitted constants as an existence proof that
a single family can cover the dimension range, not as values anyone should
reuse.

\emph{Many-Well-32.} We do not reach a competitive number on this target and
have diagnosed but not fixed the cause (\S\ref{subsec:gmm40}). Until that is
closed, the GMM-40 result should be read as one favourable external
target rather than as a general claim about Boltzmann sampling.

\emph{No pretrained model.} This is the largest gap. \S\ref{sec:02_setting}
observes that if the base process is a pretrained diffusion or flow model then
$f$ is its learned drift and nothing in the development changes, and we believe
that, but we never test it: every experiment here uses analytic base processes
--- Gaussian-mixture interpolants and driftless Brownian motion --- with
residual MLPs at gradient batch size $4$ on synthetic tilts. Whether the
conditioning differences we measure survive at the scale and with the learned
drifts of an actual reward-guided diffusion is untested, and a reader should
weight the results accordingly.

\emph{Scope of the benchmark.} The dimensional-scaling study fixes headroom at
$6$ nats and caps network width at $256$ for $d\ge8$. Both the crossover
dimension between on- and off-policy training and the size of the loss effect
plausibly move with these, and we have not varied them.

\emph{What the loss does not do.} The Spence loss improves conditioning. It
does not make rare high-value targets easier to find, which is the separate
problem of \S\ref{sec:06_sparsity_onpolicy} and the reason the two halves of
this paper are both necessary.

\paragraph{Outlook.}
The gradient-weight view suggests its own next question. Once the design space
is a single positive function $w$ on $(0,\infty)$, the weight can be chosen for
things other than tail behaviour --- robustness to target noise, or a
prescribed influence function --- and \eqref{eq:req-over}--\eqref{eq:req-under}
pin down only the asymptotics, leaving a family of admissible weights that we
have not explored.


\bibliographystyle{plainnat}
\bibliography{refs}

\appendix
% !TEX root = ../main.tex
\section{The Spence loss: derivation and implementation}
\label{sec:A_spence}

% TODO: Bregman potential F, dilogarithm identities, population-minimizer
% proof, branch-wise stable gradient, gamma-like uncertainty interpretation.

% !TEX root = ../main.tex
\section{On-policy method zoo}
\label{sec:B_onpolicy_zoo}

% TODO

% !TEX root = ../main.tex
\section{Experimental details}
\label{sec:C_experimental_details}

% TODO

% !TEX root = ../main.tex
\section{Additional proofs}
\label{sec:D_proofs}

% TODO


\end{document}
